%! AP Statistics — Unit 5, Lesson 5.1 — Warm-Up (Answer Key)
\documentclass[10pt]{article}
\usepackage{apstats-article}
\usepackage{apstats-key}   % pulls in -boxes; do NOT also load -boxes
\newcommand{\TallMath}[1]{$\displaystyle #1\rule[-1.4em]{0pt}{3.2em}$}

\begin{document}

\pageheader{Unit 5, Lesson 5.1}{Warm-Up --- Answer Key}
\namedateperiod

\vspace{0.15in}

\begin{notesbox}{Spiral Review --- Answer Key}
\textbf{1.} In Unit 1 you learned about statistics and parameters. Match each value to
the correct label (\emph{statistic} or \emph{parameter}).

\vspace{0.1in}
\begin{enumerate}[label=\alph*.]
  \item The mean SAT score of \emph{all} U.S. high school seniors last year.
        \hfill \ans{parameter}
  \item The mean SAT score of the 200 seniors surveyed at one school.
        \hfill \ans{statistic}
  \item The proportion of adults in a random sample of 1,000 who own a pet.
        \hfill \ans{statistic}
  \item The true proportion of all U.S. adults who own a pet.
        \hfill \ans{parameter}
\end{enumerate}

\vspace{0.15in}
\textbf{2.} Two students each randomly select 20 classmates and ask how many hours of
sleep they got last night.
\begin{itemize}
  \item Student A gets a sample mean of 7.2 hours.
  \item Student B gets a sample mean of 6.8 hours.
\end{itemize}

Does the difference in these sample means prove that one of them made an error?
Explain. \ansline{No. Different random samples from the same population will produce
different sample means just by chance. The 0.4-hour difference is an example of
sampling variability, not an error.}
\end{notesbox}

\begin{teachernote}
Watch for students who say ``one of them is wrong'' --- redirect to the idea that
random sampling naturally produces variation. This misconception is exactly what the
unit is designed to correct.
\end{teachernote}

\end{document}
